%% D-P2-007: definitional CNF for the active primitive CPL grammar.
\section{Definitional CNF and the Tseitin transformation}
\label{sec:p02-tseitin-transformation}

The normal forms of Section~\ref{sec:prop-syntax-semantics-04} preserve truth
under each valuation.  For satisfiability search, definitional CNF instead
uses additional variables to record subformula values.  Harrison's
author-hosted presentation gives biconditional definitions together with an
asserted output variable \cite[slide~12]{Harrison2003PropositionalSlides}.
The construction below fixes every primitive case and its occurrence and
encoding conventions.  The conjunction clause shape and the root assertion
also occur in the Part~\ref{part:p12} example in
Section~\ref{sec:automated-reasoning-02}.

Use the actual grammar \(\mathsf{Fm}_{\mathsf{CPL}}\), with primitives
\(\bottom,\neg,\land,\lor,\to\), and truth values
\(\mathbb{B}=\{\mathbf{F},\mathbf{T}\}\).  Defined connectives are expanded before applying
this construction.  A \emph{valuation on} a finite set \(W\subseteq\mathsf{Prop}\)
means a function \(v:W\to\mathbb{B}\).  The same recursive truth clauses define
\(\widehat v\) on formulas whose variables lie in \(W\).
Lemma~\ref{lem:prop-coincidence} identifies these values with those under any
total valuation extending \(v\); such an extension exists by assigning
\(\mathbf{F}\) outside \(W\).  Clause satisfaction uses
\(\operatorname{form}_{\mathsf{CPL}}\) from the preceding normal-form section.
For a finite clause set \(S\), write \(v\entails S\) when \(v\) satisfies
every clause in \(S\).

\begin{hbdefinition}{def:p02-tseitin-construction}{Occurrence-based definitional CNF}
Fix \(\varphi\in\mathsf{Fm}_{\mathsf{CPL}}\).  Let \(V\) be exactly its finite set of
occurring variables, and use \(\operatorname{Pos}(\varphi)\) and
\(\varphi_a\) from Definition~\ref{def:p02-selected-replacement}.
The root is \(\epsilon\); the children of a position \(a\) are \(a1\)
and, in the binary case, \(a2\).  Put
\[
 N=|\operatorname{Pos}(\varphi)|,\qquad
 G_\varphi=\{a\in\operatorname{Pos}(\varphi):\varphi_a\notin\mathsf{Prop}\},
 \qquad M=|G_\varphi|.
\]
Choose distinct variables \(z_a\in\mathsf{Prop}\setminus V\), one for each
\(a\in G_\varphi\), and let \(Z_\varphi=\{z_a:a\in G_\varphi\}\).
One deterministic policy traverses the tree in preorder (root, left subtree,
right subtree) and assigns the first \(M\) unused variables in the indexed
enumeration \(p_0,p_1,\ldots\).  Equal subformula strings at different
nonvariable positions receive different labels.  The notation \(z_a\) is an
alias for its chosen variable \(p_n\), not an extra kind of variable token.

Each position has a variable representative
\[
 r_a=\begin{cases}
       p_n,&\varphi_a=p_n,\\
       z_a,&a\in G_\varphi.
      \end{cases}
\]
Thus original variable leaves are reused, and a \(\bottom\) occurrence is
labelled.  Emit the following clauses at each node.  In the nonvariable rows
write \(z=z_a\), \(x=r_{a1}\) when there is a child, and
\(y=r_{a2}\) when there is a second child.  These are variables, so their
negations are literals; \(x\) and \(y\) are allowed to coincide.
\[
\begin{array}{c|l}
 \varphi_a&\text{clauses emitted at }a\\ \hline
 p_n&\text{none}\\
 \bottom&\{\neg z\}\\
 \neg\alpha&\{\neg z,\neg x\},\ \{z,x\}\\
 (\alpha\land\beta)&
   \{\neg z,x\},\ \{\neg z,y\},\ \{z,\neg x,\neg y\}\\
 (\alpha\lor\beta)&
   \{\neg z,x,y\},\ \{z,\neg x\},\ \{z,\neg y\}\\
 (\alpha\to\beta)&
   \{\neg z,\neg x,y\},\ \{z,x\},\ \{z,\neg y\}.
\end{array}
\]
Let \(\mathcal D_\varphi\) be the union of these finite clause sets, and put
\[
 S_\varphi=\mathcal D_\varphi\cup\bigl\{\{r_\epsilon\}\bigr\}.
\]
The last clause is the \emph{root clause}.  It asserts the representative of
the complete input, including when that input is just a variable.

For an exact CNF formula string, emit node blocks in preorder and the clauses
within each row in their displayed order, then append the root clause.
Call this list \(C_1,\ldots,C_m\), retaining any duplicate emissions, and set
\[
 T_\varphi=\bigwedge_{i=1}^{m}\operatorname{form}_{\mathsf{CPL}}(C_i)
\]
with the inherited right-associated convention.  Each individual clause's
literals are sorted by the inherited literal code when
\(\operatorname{form}_{\mathsf{CPL}}\) is applied.  Thus the emitted list determines
one formula string, while \(S_\varphi\) is a set of clauses.
\index{Tseitin transformation}\index{CNF!definitional}
\index{root clause}\index{auxiliary variable}
\end{hbdefinition}

\begin{hbtheorem}{thm:p02-tseitin}{proved=annex:p2.prop.tseitin}{Definitional extension and equisatisfiability}
For every \(\varphi\in\mathsf{Fm}_{\mathsf{CPL}}\), use the variable sets and clause
construction of Definition~\ref{def:p02-tseitin-construction}.
For every valuation \(v:V\to\mathbb{B}\), there is exactly one extension
\(w:V\cup Z_\varphi\to\mathbb{B}\) satisfying \(\mathcal D_\varphi\).
It satisfies
\[
 w(r_a)=\widehat v(\varphi_a)
 \qquad\text{for every }a\in\operatorname{Pos}(\varphi).
\]
Consequently,
\[
 v\entails\varphi
 \quad\Longleftrightarrow\quad
 \exists w:V\cup Z_\varphi\to\mathbb{B}\;
   [\,w|_V=v\ \text{and}\ w\entails S_\varphi\,].
\]
When the right side holds, its extension is unique on \(V\cup Z_\varphi\).
The clause set \(S_\varphi\) and the CNF formula \(T_\varphi\) are
equisatisfiable with \(\varphi\).  Under total valuations on \(\mathsf{Prop}\),
\(\{T_\varphi\}\entails\varphi\), but the converse need not hold.
\end{hbtheorem}

The existence quantifier in this theorem belongs to the metatheory, not to
the CPL formula grammar.  The input assignment is fixed on \(V\), not on
all of \(\mathsf{Prop}\): a total global valuation already assigns the chosen fresh
symbols.  Total extensions outside \(V\cup Z_\varphi\) remain unrestricted
and are not claimed unique.  Uniqueness uses both directions of every gate
definition; it is not asserted for encodings obtained by dropping selected
gate clauses.

\begin{hbtheorem}{thm:p02-tseitin-size}{proved=annex:p2.prop.tseitin-size}{Structural size and an explicit encoding bound}
For every input \(\varphi\in\mathsf{Fm}_{\mathsf{CPL}}\), with \(N,M\) as in
Definition~\ref{def:p02-tseitin-construction},
the construction uses exactly \(M\leq N\) auxiliary variables, emits at most
\(3M+1\) clauses, and emits at most \(7M+1\) literal occurrences.  Every
clause has at most three literals.  The emitted list and \(T_\varphi\) have
\(O(N)\) size when each variable name is one token.  Given the parsed tree
and the fresh names, clause emission takes \(O(N)\) elementary operations
with unit-cost variable references and name comparisons.

For a bit encoding in which variable indices are written in binary and the
connectives and separators have fixed codes, let \(L\) be the length of that
primitive input.
The deterministic policy can choose every auxiliary index below \(N\).
The output bit length is
\[
 O\bigl(L+N\log(N+1)\bigr),
\]
and the construction is effective in polynomial time in \(L\).
No linear bit-time bound is asserted.  These bounds measure the primitive
formula tree after expansion of defined connectives; they do not measure an
unexpanded abbreviation or the number of distinct subformula strings.
\end{hbtheorem}

\begin{hbexample}{def:p02-tseitin-worked}{A conjunction, negation and implication}
Take \(\varphi=((P\land Q)\to\neg P)\), with distinct \(P,Q\).
The nonvariable positions are \(\epsilon,1,2\), represented by fresh
\(z_\epsilon,z_1,z_2\).  The emitted blocks, in order, are
\[
\begin{aligned}
 \epsilon:&\quad
  \{\neg z_\epsilon,\neg z_1,z_2\},
  \{z_\epsilon,z_1\},\ \{z_\epsilon,\neg z_2\},\\
 1:&\quad
  \{\neg z_1,P\},\ \{\neg z_1,Q\},\ \{z_1,\neg P,\neg Q\},\\
 2:&\quad \{\neg z_2,\neg P\},\ \{z_2,P\}.
\end{aligned}
\]
Append \(\{z_\epsilon\}\), giving nine emitted clauses.
If \(v(P)=\mathbf{F}\), \(v(Q)=\mathbf{T}\), the defining extension gives
\(w(z_1)=\mathbf{F}\), \(w(z_2)=\mathbf{T}\), \(w(z_\epsilon)=\mathbf{T}\),
and satisfies the output.  If both input variables are true, the defining
extension instead has values \(\mathbf{T},\mathbf{F},\mathbf{F}\) at these three
labels: it satisfies every gate block but fails the root clause, just as the
input formula is false.  The three clauses at position \(1\) have the same
conjunction-gate shape as the Part~\ref{part:p12} example; they do not assert the root.
\end{hbexample}

\begin{hbexample}{def:p02-tseitin-root-boundary}{Why the root clause and the assignment domains matter}
For \(\varphi=\bottom\), the input variable set is empty and
\[
 \mathcal D_\varphi=\bigl\{\{\neg z_\epsilon\}\bigr\},\qquad
 S_\varphi=\bigl\{\{\neg z_\epsilon\},\{z_\epsilon\}\bigr\}.
\]
The definitions alone are satisfiable, while the input and the rooted output
are unsatisfiable.  Dropping the root therefore loses equisatisfiability.

For \(\varphi=\neg P\), a valuation on \(\{P,z_\epsilon\}\) assigning
\(\mathbf{F}\) to both variables makes \(\varphi\) true and \(T_\varphi\)
false.  Hence \(\varphi\) and \(T_\varphi\) are not generally logically
equivalent even when both are read over their common enlarged vocabulary.
The extension theorem instead fixes \(P\) and chooses the required fresh
value \(w(z_\epsilon)=\mathbf{T}\).
\end{hbexample}
